\input amstex
\documentstyle{amsppt}
\nologo
\magnification=\magstep1
\parindent=15pt
\pageno=1
\vsize=23truecm
\hsize=16truecm
\baselineskip=19pt plus 3pt minus 3pt
\document
\topspace{0.5cm}
\baselineskip=19pt plus 3pt minus 3pt
\centerline{{\bf A group-theoretic characterization of the space obtained by omitting}}
\centerline{{\bf the coordinate hyperplanes from the complex Euclidean space}}
\vskip 1cm
\noindent\hskip 7.2cm Akio Kodama \;\;(Kanazawa Univ.)\par
\noindent\hskip 7.2cm Satoru Shimizu\;\; (Tohoku Univ.)\par
\vskip 0.8cm
In the study of the holomorphic automorphism group
$\text{\rm Aut}(M)$ of a complex manifold $M$,
it seems to be natural to direct our attention to
not only the abstract group structure of $\text{\rm Aut}(M)$
but also the topological group structure of
$\text{\rm Aut}(M)$ equipped with the compact-open topology.
In fact, a well-known theorem of H. Cartan says that
the topological group given as 
the holomorphic automorphism group of
a bounded domain in $\bold C^n$
has the structure of a Lie group,
and this result enables us to make various kinds of
detailed studies of bounded domains in $\bold C^n$.
On the other hand, in contrast to the case of bounded domains,
the holomorphic automorphism group
$\text{\rm Aut}
 (\bold C^k
 \times
 (\bold C^{\ast})^{\ell})$
of the unbounded domain
$\bold C^k \times (\bold C^{\ast})^{\ell}$
is terribly big when $k + \ell \ge 2$,
and can not have the structure of a Lie group.
But, by looking at topological subgroups of 
$\text{\rm Aut}
 (\bold C^k
 \times
 (\bold C^{\ast})^{\ell})$
with Lie group structures,
we can find a lead to apply the Lie group theory to
the investigation of the problems related to
the structure of
$\text{\rm Aut}
 (\bold C^k
 \times
 (\bold C^{\ast})^{\ell})$.
In fact, we can prove the following theorem with the aid of the theory of Reinhardt domains developed in Shimizu [7], [8], and the fundamental result on torus actions on complex manifolds due to Barrett-Bedford-Dadok [2](cf. [6], [3]):\par
\vskip 0.4cm
{\bf Main Theorem.}
{\it Let $M$ be a connected complex manifold of dimension $n$
that is holomorphically separable and
admits a smooth envelope of holomorphy.
Assume that
$\text{\rm Aut}(M)$ is isomorphic to
$\text{\rm Aut}
 (\bold C^k
 \times
 (\bold C^{\ast})^{n - k})$
as topological groups.
Then $M$ itself is biholomorphically
equivalent to
$\bold C^k
 \times
 (\bold C^{\ast})^{n - k}$.}
\par
\vskip 0.4cm
{\bf Corollary.}
{\it If two pairs
$(k, \ell)$
and
$(k^{\prime}, \ell^{\prime})$
of nonnegative integers
do not coincide,
then the topological groups
$\text{\rm Aut}
 (\bold C^k
 \times
 (\bold C^{\ast})^{\ell})$
and
$\text{\rm Aut}
 (\bold C^{k^{\prime}}
 \times
 (\bold C^{\ast})^{\ell^{\prime}})$
are not isomorphic.}
\par
\vskip 0.4cm
It should be remarked that, as shown in Ahern and Rudin [1], the groups $\text{\rm Aut}({\bold C}^n)$ and $\text{\rm Aut}({\bold C}^m)$ are isomorphic as abstract groups precisely when $n=m$. Also, our method can be applied to the study of unitary group actions on complex manifolds. The following theorem and its corollary give a different approach from Kaup [5], Isaev-Kruzhilin [4] to the study of $U(n)$-actions on a complex manifold of dimension $n$:
\par
\vskip 0.4cm
{\bf Theorem.}
{\it Let $M$ be a connected Stein manifold
of dimension $n \ge 2$
and let $K$ be a compact group given as
the direct product
$K = U(n_{1}) \times \cdots \times U(n_{s})$
of unitary groups,
where each $U(n_{j})$ is the unitary group
of degree $n_{j} > 1$ and
$\sum_{j = 1}^{s} n_{j} = n$.
Suppose that
an injective continuous group homomorphism
$\rho$ of $K$ into $\text{\rm Aut}(M)$
is given.
Then there exists a biholomorphic mapping $F$
of $M$ onto a Reinhardt domain $D$ in $\bold C^n$
such that $K$ is a subgroup of $\text{\rm Aut}(D)$
and $F \rho(K) F^{-1} = K$.}
\par
\vskip 0.4cm
{\bf Corollary.}
{\it In the above theorem,
if $K = U(n)$, or $M$ admits an effective
$U(n)$-action, then
$M$ is biholomorphically equivalent to
either the unit ball in $\bold C^n$
or $\bold C^n$ itself.}
\par
\vskip 0.5cm
\centerline{{\bf References}}
\vskip 0.7cm
\item{[1]} P. Ahern and W. Rudin:  {\it Periodic Automorphisms of ${\bold C}^n$},  Indiana Univ\. Math\. J. {\bf 44} (1995), 287--303.
\item{[2]} D. E. Barrett, E. Bedford, and J. Dadok:  {\it $T^n$-actions on holomorphically separable complex manifolds},  Math\. Z. {\bf 202} (1989), 65--82.
\item{[3]} A. V. Isaev:  {\it Characterization of $\bold C^n$ by its automorphism group},  Proc\. Steklov Inst\. Math\. 2001, no\. 4 (235), 103--106.
\item{[4]} A. V. Isaev and N. G. Kruzhilin:  {\it Effective actions of the unitary group on complex manifolds},  Canad\. J\. Math. {\bf 54} (2002), 1254--1279.
\item{[5]} W. Kaup:  {\it Reelle Transformationsgruppen und invariante Metriken auf komplexen R\"aumen},  Invent.\ Math.\ {\bf 3} (1967), 43--70.
\item{[6]} A. Kodama and S. Shimizu:  {\it A group-theoretic characterization o
the space obtained by omitting the coordinate hyperplanes
from the complex Euclidean space},  Osaka J\. Math. {\bf 41} (2004), 1--11.
\item{[7]} S. Shimizu:  {\it Automorphisms and equivalence of bounded Reinhardt domains not containing the origin},  Tohoku Math.\ J. {\bf 40} (1988), 119--152
\item{[8]} S. Shimizu:  {\it Automorphisms of bounded Reinhardt domains},  Japan.\ J. Math.\ {\bf 15} (1989), 385--414.
\enddocument
